% ============================================================================
% 第 3 章 系统设计与实现（合并自原 ch03 模块设计 + ch04 系统实现）
% ============================================================================

\section{系统设计与实现}

第二章给出了伏羲的整体架构与通讯模型形式化。本章把这套架构落到 Rust 实现，按模块依赖自下而上的顺序展开七个核心单元的设计与实现：模块布局、事件总线、A2A 协议、编排与任务派发、执行适配器与三层沙箱、跨节点通讯、观测与横向支撑。每节遵循同一主线：先讲外部约束与关键数据结构、再用一段 Rust 代码片段或公式呈现工程上不寻常的部分、最后追加关键设计取舍。本章保留的代码片段均节选自主仓 \texttt{/Users/e0\_7/fuxi}（HEAD: \texttt{04eb34c}）下对应文件，为版面计做了行宽折断与无关字段省略，未修改控制流。

为给整章提供一张索引图，图~\ref{fig:execution-loop} 把伏羲一次完整任务执行的端到端闭环画在一张图上：用户从 IM/PWA/REPL 输入指令（步骤 1），玄女接到 \texttt{intervene} 后向 \texttt{Shelf} 申领空闲门客（步骤 2），通过 A2A \texttt{tasks/sendSubscribe} 把任务发到门客的 JSON-RPC + SSE 端点（步骤 3），门客在 L2 git worktree 沙箱中执行（步骤 4）；执行过程中产生的所有事件向 EventBus publish（步骤 5），由 EventBus 同时走 broadcast 实时扇出（步骤 6a）与 WAL 串行落库（步骤 6b）；观测端断线重连时通过 \texttt{replay(cursor)} 从 SQLite 拉历史补齐（步骤 7），再无缝拼接到实时流；若门客进入 \texttt{InputRequired} 中间态，事件被翻译成 \texttt{ShelfStatus::AwaitingInput}（步骤 8），玄女据此请求人工回复（步骤 9）。后续小节展开的代码与算法均服务于这同一个闭环。

\begin{figure}[H]
  \centering
  \includegraphics[width=0.96\textwidth]{drawio/fig-3-1-execution-loop.pdf}
  \caption{伏羲端到端任务执行闭环（9 步从用户提交至多端观测）}
  \label{fig:execution-loop}
\end{figure}

\subsection{模块布局与 13 crate 边界}
\label{subsec:module-layout}

伏羲在 crate 粒度上划分 13 个模块，总代码规模 82\,515 行 Rust 源（\texttt{find crates -name "*.rs" \textbar{} xargs wc -l}）。各 crate 按职责分层组织，依赖方向严格自下而上、禁止循环引用，整体如图~\ref{fig:module-dependency}。表~\ref{tab:crate-layout} 列出每个 crate 的代码规模、对外类型与主要职责。

\begin{figure}[H]
  \centering
  \includegraphics[width=0.85\textwidth]{drawio/fig-3-2-module-dependency.pdf}
  \caption{伏羲模块依赖关系（13 crate，5 层架构）}
  \label{fig:module-dependency}
\end{figure}

\begin{table}[H]
\centering
\caption{伏羲 workspace 13 个 crate 的边界与代码规模}
\label{tab:crate-layout}
\footnotesize
\begin{tabularx}{\textwidth}{p{2.6cm} r p{3.4cm} X}
\toprule
\textbf{Crate} & \textbf{LOC} & \textbf{对外类型/trait} & \textbf{主要职责} \\
\midrule
fuxi-cli & 33\,408 & 二进制入口 & REPL + daemon + dist controller + 钩子注入 \\
fuxi-im & 16\,149 & axum app, REST/WS handlers & IM 后端：任务/对话/intervene 端点 \\
fuxi-orchestrator & 10\,220 & \texttt{Fuxi}, \texttt{Shelf}, \texttt{RecallSink} & 编排中枢：注册表 + dispatch + republish \\
fuxi-agent-cc & 3\,900 & \texttt{CcAgent} & Claude Code 适配器（WS 反连） \\
fuxi-workspace & 3\,043 & \texttt{Git\-Worktree\-Workspace} & git worktree 隔离 + L1/L2/L3 沙箱 \\
fuxi-memory & 2\,643 & \texttt{OracleStore}, \texttt{ExtractorSpawner} & 长期记忆：事实/技能/身份 \\
fuxi-agent-codex & 2\,197 & \texttt{CodexAgent} & Codex 适配器（spawn-per-dispatch） \\
fuxi-scheduler & 2\,146 & \texttt{Keeper}, \texttt{Watcher} & cron/once/webhook/fs 触发器 \\
fuxi-events & 2\,135 & \texttt{EventBus}, \texttt{EventStore} & 实时分发 + WAL 持久化 + 历史回放 \\
fuxi-firehose & 2\,135 & \texttt{Hub}, axum routes & 事件流四端输出（WS/SSE/REST/TUI） \\
fuxi-core & 2\,102 & \texttt{Agent}, \texttt{Workspace}, \texttt{Event}, \texttt{Task} & 核心 trait 与值对象，零外部 IO \\
fuxi-a2a & 1\,321 & \texttt{wire::*}, \texttt{Server}, \texttt{Client} & A2A wire + axum 服务器 + reqwest 客户端 \\
fuxi-skills & 1\,116 & \texttt{RoleLoader}, \texttt{SkillLedger} & 角色玉牒加载与招贤流程 \\
\midrule
\textbf{总计} & \textbf{82\,515} & & \\
\bottomrule
\end{tabularx}
\end{table}

\texttt{fuxi-cli} 占 workspace 近 44\%——这反映大型 Rust 项目的工程现实：协议、抽象与基础设施虽然概念上是「重头戏」，但把它们装配成一个能跑的产品需要大量胶水代码。模块边界设计的隐含原则是把会高频变化的「策略」与会长期稳定的「机制」分开放在不同 crate：\texttt{fuxi-core} 中的 trait 与 \texttt{EventKind} 枚举属于稳定层，\texttt{fuxi-orchestrator} 内部调度策略与 \texttt{fuxi-agent-cc} 内事件翻译规则属于策略层，迭代频率高但不影响其他 crate。

依赖图中存在一处需绕过环路的需求：\texttt{fuxi-memory} 在召回结果回传时需要往 orchestrator 持有的事件总线 publish 事件，但若让 memory 反向依赖 orchestrator 会形成循环。本设计采用「trait 在低层、impl 在高层」的反向依赖：在 orchestrator 中定义 \texttt{RecallSink} trait，由顶层 \texttt{fuxi-cli} 提供实现并在 daemon 启动时注入；类似地有 \texttt{ExtractorSpawner}、\texttt{DistEnqueuer}、\texttt{NodeLoadProvider} 三个钩子，统一定义在低层 crate、由 \texttt{fuxi-cli} 注入实现。这是规避 Rust workspace 循环依赖的标准 Robert Martin 依赖反转范式。

\subsection{事件总线模块}
\label{subsec:eventbus}

事件总线由 \texttt{fuxi-events} 实现，是公理「真实时不轮询」与「SQLite 是单一真相源」的工程落地。该模块同时承担三件事：实时分发、持久化、历史回放。一个 \texttt{Event} 由 \texttt{EventMeta}（携带通用元数据）与 \texttt{EventKind}（带 50+ 变体的 \texttt{\#[serde(tag = "type")]} 标签联合）构成。\texttt{tag = "type"} 的关键价值是反序列化无歧义，且新增变体时所有 \texttt{match ev.kind} 处都会编译失败强制更新所有消费路径——这是 Rust enum 与 \texttt{serde} 在类型安全上的典型应用。

\textbf{双路径架构。} \texttt{EventBus::publish()} 在每次调用时同时把事件投递到两条独立路径：第一条是 Tokio \texttt{broadcast} channel\parencite{tokio2024broadcast} 零拷贝扇出给所有订阅者实现实时推送；第二条是异步 \texttt{mpsc} 队列由后台 writer task 串行落库到 SQLite，启用 WAL 模式\parencite{sqlite2024wal} 让读写并发不互斥。两条路径各走最优形态，避免「单条通道既当广播又当持久化」的拥塞反模式。

\textbf{非阻塞发布与 lag 哨兵。} 发布路径核心约束是「调用方永不阻塞」与「事件不丢」两个看似矛盾的约束，通过「\texttt{try\_send} + 后台转交 + 哨兵告警」三步组合解决，见代码片段~\ref{lst:eventbus-publish-snippet}。

\begin{lstlisting}[language=Rust, caption=EventBus::publish 非阻塞实现节选（fuxi-events/src/bus.rs）, label=lst:eventbus-publish-snippet]
pub fn publish(&self, ev: Event) -> Result<()> {
    // 路径 1: broadcast 零拷贝扇出, 不影响 mpsc
    let _ = self.inner.broadcast_tx.send(ev.clone());
    // 路径 2: try_send 入 mpsc（非阻塞）
    match self.inner.writer_tx.try_send(ev) {
        Ok(_) => { self.inner.writer_pending.fetch_add(1, Ordering::Relaxed); }
        Err(TrySendError::Full(dropped)) => {
            // 拥塞: 后台 spawn 等慢车位, 当前 publish 立即返回
            let pending = self.inner.writer_pending.load(Ordering::Relaxed);
            if pending >= self.inner.cfg.lag_threshold {
                self.maybe_emit_lag_sentinel(pending);    // 哨兵告警
            }
            let tx = self.inner.writer_tx.clone();
            tokio::spawn(async move { let _ = tx.send(dropped).await; });
        }
        Err(TrySendError::Closed(_)) => return Err(Error::WriterClosed),
    }
    Ok(())
}
\end{lstlisting}

\texttt{maybe\_emit\_lag\_sentinel} 通过同一个 broadcast 路径发出一条 \texttt{Custom \{ "event\_store\_lagged", pending \}} 哨兵事件，订阅者据此感知到底层 writer 已经堆了多少未落库事件。哨兵自身有去重逻辑，不会反向加剧拥塞。该设计与 Tower 在 Rust 服务栈中的 \texttt{tower::load} 背压抽象\parencite{tower_backpressure_2024} 方向一致——都把背压视为一等公民信号而非异常状态；NATS\parencite{nats_2024} 与 Kafka 同样把订阅者滞后纳入显式监控指标。替代方案有二：直接丢消息违反「事件不丢」公理；让 \texttt{publish()} 阻塞等待 mpsc 容量会让上层编排在偶发拥塞下卡死。两种简单方案均不可取。

\textbf{回放与因果序。} \texttt{EventBus::replay(cursor, live\_tail)} 接受 \texttt{ReplayCursor}（\texttt{Beginning}/\texttt{FromId}/\texttt{FromTime}），先从 SQLite 拉历史流，若 \texttt{live\_tail = true} 则在历史末尾无缝拼接 broadcast 实时流。实现上必须先订阅 broadcast 再开始拉历史——反过来会留出窗口，期间产生的事件既未入历史快照也未被订阅者收到。事件因果序遵循 Lamport 逻辑时钟思路\parencite{lamport1978time}：伏羲在每条事件 \texttt{meta} 字段记录 Agent ID、Task ID、墙钟时间戳与 SQLite \texttt{rowid}（单调递增），订阅者重建因果链时用 \texttt{rowid} 排序，避免引入向量时钟所需的时钟同步基础设施。

\textbf{SQLite 工程细节。} 事件持久化表 \texttt{events} 包含 id/at/session/agent/task/kind\_tag/payload 七列，\texttt{kind\_tag} 独立索引使按类型过滤为 O(log n) 扫描，\texttt{payload} 存放完整 \texttt{Event} JSON 保证 schema 演进时零破坏迁移。WAL 模式高并发下偶发 \texttt{SQLITE\_BUSY}，伏羲通过 \texttt{busy\_timeout=5s} 加三次线性退避（50/100/150 ms）封装在 \texttt{EventStore::append} 内部，外层调用方无感知。\texttt{:memory:} 库每条连接独立持有内存库，\texttt{EventStore::new} 检测到时强制 \texttt{max\_connections = 1} 规避「明明 append 了却 replay 不出来」的诡异行为。整个 \texttt{fuxi-events} crate 共 1673 行（store.rs 1024 + bus.rs 649），是从外观简单的 publish/subscribe/replay 接口背后挖出工程取舍最多的模块之一。

\subsection{A2A 协议模块（论文核心贡献）}
\label{subsec:a2a}

通讯协议模块由 \texttt{fuxi-a2a} 实现，是本论文最重要的工程贡献之一。Rust 生态已有 \texttt{a2a-rs}（A2A 项目官方 SDK）以及 \texttt{a2a-client}、\texttt{a2a-types} 等第三方 crate\parencite{a2aproject2025a2a}，本工作从零实现 wire 层、JSON-RPC 信封、SSE 流抽象、axum 服务器与 reqwest 客户端，共计 1063 行核心代码加 258 行端到端测试，与 \texttt{fuxi-events} 和 \texttt{fuxi-orchestrator} 同源演进。

\textbf{协议层次划分。} A2A 协议解决「Agent 与 Agent 之间的协作契约」，MCP\parencite{anthropic2024mcp} 解决「模型与外部工具的调用规范」，二者面向不同语义层次不互斥。伏羲内部门客之间使用 A2A 传递任务与结果，门客调用外部工具沿用其底层 CLI 自身的工具机制（cc 与 codex 都遵循 MCP 风格的工具声明）。这一分工与 SWE-agent\parencite{yang2024sweagent} 的 Agent-Computer Interface 思想一致。替代方案曾考虑把所有协议统一到 MCP 让门客之间通过虚拟工具相互调用，被否决的原因是 MCP 工具语义偏「无状态、单次调用」而 Agent 协作需要「状态机、流式响应、人工介入点」。这里需进一步区分：A2A 描述「Agent 之间应当传递什么」，事件总线描述「系统内部已经发生了什么」——前者是动作语义，后者是事实语义，伏羲明确划清这条边界让 A2A 信封只承担最小契约。

\textbf{Crate 模块拆分。} \texttt{crates/fuxi-a2a} 拆分 7 个文件：\texttt{wire}（458 行）、\texttt{jsonrpc}（77 行）、\texttt{sse}（45 行）、\texttt{server}（196 行）、\texttt{client}（203 行）、\texttt{error}（52 行）、\texttt{lib.rs}（32 行），加端到端测试 \texttt{tests/roundtrip.rs} 258 行，总计 1321 行（文档注释占 25\%）。本实现采用早期 A2A JSON-RPC binding，沿用方法命名 \texttt{agent/getCard}、\texttt{tasks/send}、\texttt{tasks/sendSubscribe}、\texttt{tasks/get}、\texttt{tasks/cancel}（当前官方规范已切到 PascalCase 方法集 \texttt{SendMessage}、\texttt{GetTask} 等，本工作未跟进改造，亦未与 \texttt{a2a-rs} 做完整 wire-level 互通验证；不主张「v1.0 兼容」）。

\textbf{wire 层标签联合。} 核心数据结构包括 \texttt{AgentCard}、\texttt{Task}、\texttt{Message}、\texttt{Part}、\texttt{Artifact}。所有结构使用 \texttt{serde} 派生序列化，并通过 \texttt{\#[serde(rename\_all = "kebab-case")]} 与 \texttt{\#[serde(tag = "type")]} 精确控制 wire 上 JSON 形态。代码~\ref{lst:a2a-part-tagged} 是 \texttt{Part} 的写法。

\begin{lstlisting}[language=Rust, caption=Part 类型的 serde 标签联合实现, label=lst:a2a-part-tagged]
#[derive(Debug, Clone, Serialize, Deserialize)]
#[serde(tag = "type", rename_all = "kebab-case")]
pub enum Part {
    Text { text: String },
    Data { data: serde_json::Value },
    File { file: FileContent },
}
\end{lstlisting}

使用 \texttt{tag = "type"} 而非 \texttt{untagged} 是经过权衡的——\texttt{untagged} 让 serde 按字段形态依次尝试匹配，在 \texttt{Text} 与 \texttt{Data} 字段近似时会产生反序列化优先级歧义；显式标签强制每个 Part 在 wire 上携带 \texttt{"type": "text"} 类判别字段，反序列化结果稳定可预测。\texttt{fuxi-core::AgentCard}（带 \texttt{id}、\texttt{status} 运行时字段）与 \texttt{fuxi-a2a::wire::AgentCard}（协议对外视图）之间的边界转换通过 \texttt{From} trait 实现：丢弃内部字段、把 \texttt{system\_prompt} 首行截 200 字作 \texttt{description}、把 \texttt{tags} 平铺成 A2A \texttt{skills}。本设计禁止在任何业务路径中手抄字段——所有边界转换必须走 \texttt{From} trait，由编译器保证两种 \texttt{AgentCard} 同步演进。

\textbf{InputRequired：A2A 规范语义在伏羲平台的映射。} A2A 当前规范定义的 \texttt{TaskState} 包含 \texttt{Submitted}、\texttt{Working} 两个进行态，\texttt{Completed}、\texttt{Failed}、\texttt{Canceled} 三个终态，以及 \texttt{InputRequired}、\texttt{AuthRequired} 两个 interrupted state（非终态、可恢复回 \texttt{Working}）和 \texttt{Rejected} 这一拒绝执行终态。伏羲提前实现 \texttt{InputRequired}（对应规范后续正式纳入的 ProtoJSON 枚举 \texttt{TASK\_STATE\_INPUT\_REQUIRED}）；\texttt{AuthRequired} 与 \texttt{Rejected} 当前实现暂未覆盖（详见 §5.2）。代码~\ref{lst:taskstate} 是 fuxi-a2a 当前枚举。

\begin{lstlisting}[language=Rust, caption=fuxi-a2a TaskState 枚举（提前实现 InputRequired 中间态）, label=lst:taskstate]
#[derive(Debug, Clone, Copy, PartialEq, Eq, Serialize, Deserialize)]
#[serde(rename_all = "kebab-case")]
pub enum TaskState {
    Submitted,      // 刚入队
    Working,        // 处理中
    InputRequired,  // 需要人工输入（早期 draft 未列、fuxi 提前实现）
    Completed, Failed, Canceled,  // 三个终态
}
\end{lstlisting}

这一平台映射的工程价值在于把人工介入从「编排器内部隐式状态」提升为「可被协议字段承载、可被事件流观测的显式状态」。玄女订阅门客的 SSE 流时一旦看到 \texttt{state: input-required} 可立即把对话权切换回用户，等用户回复后再走 \texttt{tasks/send} 把答复回传给同一 Task ID。在伏羲实际使用中这一机制覆盖「门客对需求理解有歧义」「门客对工具调用授权犹豫」「门客遇到外部依赖问题」三类场景，全部通过同一状态语义统一处理。\texttt{InputRequired} 与终态的关键区别是：终态不可恢复，事件流关闭；\texttt{InputRequired} 是中间态，事件流保持打开，等用户回复后转回 \texttt{Working}。这条「中间态可恢复」让多轮对话与多回合协作能在同一 Task 上下文内完成。进一步地 \texttt{InputRequired} 与编排层 \texttt{ShelfStatus::AwaitingInput} 一一对应——使整体调度器（包括跨节点 \texttt{NodeLoadProvider}）能将「在等用户」的门客负载与「在工作」的负载区分开。

\textbf{单端点 HTTP+SSE 与帧解析。} 伏羲只暴露一个 HTTP 端点 \texttt{POST /a2a}，所有方法通过 JSON-RPC 2.0 信封 \texttt{method} 字段分发；\texttt{tasks/sendSubscribe} 是流式方法，对它的响应被升级为 \texttt{text/event-stream}\parencite{whatwg2024sse} 在同一 TCP 连接上回传任务进展。这一设计避免客户端管理两条独立连接的同步问题。SSE 帧的客户端解析需处理 TCP chunk 边界上的不完整帧，伏羲的 \texttt{parse\_sse\_stream} 用增量缓冲在每次收到字节后查找 \texttt{\textbackslash n\textbackslash n} 分隔符，已确证能在跨包边界场景下正确切帧（算法~\ref{lst:a2a-sse-parse}）。

\begin{lstlisting}[language=Pseudo, caption=A2A 客户端 SSE 帧解析（fuxi-a2a/src/client.rs::parse\_sse\_stream）, label=lst:a2a-sse-parse]
Input:  byte_stream
Output: 流式产出 ServerSentEventPayload
 1: buf <- ""
 2: while chunk <- byte_stream.next() do
 3:     append chunk (utf8) to buf
 4:     while buf 含 "\n\n" 边界 do
 5:         (frame, rest) <- buf.split_at("\n\n")
 6:         buf <- rest
 7:         payload <- decode_frame(frame)   // 多行 data: 拼接, : 注释跳过
 8:         yield payload
 9:     end while
10: end while
\end{lstlisting}

性能上选 \texttt{String} 增量缓冲 + \texttt{find("\textbackslash n\textbackslash n")} 朴素策略而非 \texttt{tokio\_util::codec::FramedRead} 工业级 codec——SSE 帧体积通常很小（status < 4KB、artifact < 64KB），单帧重扫开销可忽略，codec 路线代码量翻倍但收益微乎其微。\texttt{tests/roundtrip.rs} 包含三组关键测试：\texttt{roundtrip\_all\_methods} 跑完 5 个 RPC 方法的完整调用闭环；\texttt{roundtrip\_subscribe\_stream} 在 server 端故意把 SSE 帧切成 1 字节 chunks 发出验证最恶劣 chunk 边界下的正确性；\texttt{unknown\_method\_returns\_jsonrpc\_error} 验证 server 对未知方法返回 200 + \texttt{-32601} error 而非 HTTP 4xx——这条 invariant 是 JSON-RPC 规范要求。

\subsection{编排与任务派发模块}
\label{subsec:orchestrator}

编排模块由 \texttt{fuxi-orchestrator} 实现是整个系统的中枢，核心职责包括：维护门客注册表与生命周期、把任务派给合适门客、提供事件 republish 通路。该模块代码量约 8039 行。

\textbf{玄女—门客分层。} 伏羲采用「玄女顶层 Agent + 门客 worker Agent」分层结构，与 MetaGPT\parencite{hong2023metagpt}「角色分工 + 显式协作」思想一致但更强调编排层不可绕过——所有门客之间通讯必须经玄女知情，任何调度决策都发出对应事件。这条公理在 \texttt{Fuxi::shutdown\_agent} 中反复实施：所有 shutdown 路径必须比对玄女 ID，命中则 warn 并 noop 返回，唯一例外是显式 \texttt{shutdown\_xuannv\_for\_handoff}——该命名本身即作为约束，使 grep 即可识别「不要在别处调用」的意图。注册表 \texttt{Shelf} 用 \texttt{Arc<RwLock<HashMap<AgentId, ShelfEntry>>>} 实现，工作模式是「大量读 + 少量写」符合 \texttt{RwLock} 偏向读优化语义。

\textbf{Worker 选择函数。} 给定任务 $t$ 与门客集合 $\mathcal{W}$，伏羲调度函数选择满足约束的最低负载者：
\begin{equation}
\text{select}(t) = \arg\min_{w \in \mathcal{W}} \left\{ \frac{w.\text{inflight}}{w.\text{max\_concurrency}} \;\middle|\; \text{tags}(t) \subseteq \text{tags}(w),\ \text{pin}(t) \in \{\emptyset,\ w.\text{node}\} \right\}
\end{equation}
三层过滤：标签兼容性、节点钉死约束、最低饱和度。\texttt{max\_concurrency = 0} 边界归一化为 1 避免除零。这套规则与 MapReduce\parencite{dean2004mapreduce}「数据本地性 + 负载均衡」思路一致但更轻量。\texttt{pinned\_node} 优先级守卫：当用户在 IM 端显式 \texttt{@<node>} 钉死任务时 \texttt{NodeLoadProvider::auto\_pin\_from\_project} 必须立即返回 \texttt{None}——不能再悄悄改写 \texttt{pinned\_node}；调用方再加 \texttt{task.pinned\_node.is\_none()} 判断作为双保险。这是分布式自动负载均衡与用户显式意图之间的边界明确化。

\textbf{原子 claim 与 dispatch pump。} 早期版本将 \texttt{find\_idle\_by\_role} 与 \texttt{set\_status(Busy)} 拆为两步，并发 dispatch 可能 \texttt{find} 到同一只门客并分别 \texttt{set\_status}，使同一门客被双派；现在 \texttt{claim\_idle\_by\_role} 在 \texttt{write().await} 锁内完成「找空闲 + 置 Busy + 清 \texttt{idle\_since}」三步原子操作（代码~\ref{lst:claim-idle}）。第二个并发陷阱是 idle GC：另一条独立线程会扫 \texttt{idle\_since} 字段超阈值的门客并 \texttt{shutdown\_agent}，若 claim 时不清掉 \texttt{idle\_since}，已 busy 的门客仍被 GC 视为长时间空闲而误杀；显式置 \texttt{None} 才能让 GC 与 claim 不冲突。

\begin{lstlisting}[language=Rust, caption=原子按角色 claim 一个空闲门客（fuxi-orchestrator/src/registry.rs）, label=lst:claim-idle]
pub async fn claim_idle_by_role(&self, role: &str) -> Option<AgentId> {
    let mut guard = self.inner.write().await;
    let pick = guard.iter_mut()
        .find(|(_, e)| e.status == ShelfStatus::Idle && e.card.profile.role == role);
    if let Some((id, e)) = pick {
        e.status = ShelfStatus::Busy;
        e.idle_since = None;  // 占用必须清, 否则 GC 误杀
        Some(*id)
    } else { None }
}
\end{lstlisting}

\texttt{Fuxi::dispatch} 主循环见算法~\ref{lst:dispatch-main}（隐去 v2 分布式分支细节，保留控制流主干）。

\begin{lstlisting}[language=Pseudo, caption=本地任务派发主循环（fuxi-orchestrator/src/fuxi.rs::dispatch）, label=lst:dispatch-main]
Input:  agent_id, task
 1: if task.pinned_node is None and task is in a project then
 2:     task.pinned_node <- auto_pin_from_project(task)   // saturation 最低
 3: end if
 4: if task.pinned_node is Some or task.required_tags is non-empty then
 5:     emit TaskCreated; call dist_enqueuer.enqueue(task); return  // 远端接管
 6: end if
 7: agent <- shelf.get_agent(agent_id)
 8: rx <- agent.dispatch(task)              // 拿门客私有 mpsc Receiver
 9: spawn pump task:
10:     saw_terminal <- false
11:     loop:
12:         ev <- rx.recv() with timeout 50ms
13:         if ev is Some(e) then
14:             e.meta.task <- task.id; publish e to EventBus
15:             if e.kind is terminal then saw_terminal <- true
16:         elif rx is closed then break
17:         elif saw_terminal then break    // grace 超时, 安全退出
18:         end if
19:     end loop
\end{lstlisting}

第 12 行 50\,ms grace 窗口是实测教训：cc 在发出 \texttt{ResultSuccess} 后还会异步 flush 几条 \texttt{DeliverableProduced} 事件，立即 break 会让它们被丢；M2.1 修复后给 50\,ms 缓冲再退全部能 republish。\texttt{auto\_pin\_from\_project} 是 v2 跨节点版本加入的：调 \texttt{NodeLoadProvider} 取每个候选节点 \texttt{saturation = inflight / max\_concurrency} 选最低，候选全离线时返 None 退化到本地 spawn。

\textbf{intervene 退化语义与 SystemEventBridge。} 当用户在玄女对话中插入新指令时，玄女若处于 Idle 状态没有活跃 \texttt{active\_tx}，直接 \texttt{send\_message} 会让消息悬空。修复思路是 \texttt{Fuxi::intervene} 入口先查 \texttt{ShelfStatus}：Idle 退化为 \texttt{dispatch} 起一轮新 turn，Busy 把消息推入 \texttt{PendingOutbox} 等下次 turn 边界 drain。这条「Idle 退化为 dispatch」规则避免了「玄女一忙就丢消息」的常见症状。\texttt{SystemEventBridge} 是独立 tokio task：启动后立即 \texttt{bus.subscribe()} 拿 broadcast Receiver，循环读取每条事件，对 \texttt{AgentDead}/\texttt{TriggerFired}/\texttt{TaskStateChanged} 等系统事件做翻译，转成「系统提示文本」通过 \texttt{Intervener::intervene\_system\_origin} 喂给玄女。\texttt{Lagged} 错误被静默 filter 保证单次落帧不杀整条订阅链。这条桥是公理「真实时不轮询」的具体落地——玄女完全不查询任何状态，所有「知情权」都通过 EventBus 推送实现。

\subsection{执行适配器与三层沙箱}
\label{subsec:adapter-workspace}

执行层把抽象 \texttt{Agent} trait 落地到具体 CLI 进程上。伏羲实现两个适配器（\texttt{fuxi-agent-cc}、\texttt{fuxi-agent-codex}），通过 \texttt{fuxi-workspace} 三层沙箱保证多门客并发文件隔离。

\textbf{cc 与 codex 适配器差异。} \texttt{CcAgent} 包装 Claude Code CLI，采用「长连接 WS + 多轮 dispatch」：进程启动一次后通过 WebSocket 反连，多次 \texttt{dispatch} 复用同一进程。\texttt{CodexAgent} 包装 Codex CLI，采用「spawn-per-dispatch」：每次 \texttt{dispatch} 起一新进程，prompt 作为位置参数传入，进程结束即销毁。差异源自 CLI 自身能力——cc 支持 SDK URL 与会话保持，codex 是 one-shot exec 模式无 stdin 输入。\texttt{CodexAgent::send\_message} 直接返回错误，意图是强制上层把「对 idle codex 门客发消息」退化为重新 spawn 一个新进程把消息作为新 prompt 喂入——用户视角无感。

\textbf{WS 反连关键工程细节。} 见代码~\ref{lst:cc-launch}。\texttt{launch\_with\_id} 强制传入 \texttt{id} 参数——曾踩过的坑是编排层和适配器各自生成 \texttt{AgentId}，导致 \texttt{AgentSpawning} 与 \texttt{AgentReady} 事件挂在不同 ID 上订阅者无法匹配。这条新规约是「新 agent 适配器三铁律」之一：必须接收外部传入的 ID，绝不自己生成；事件 pump 无论怎么退出都要把 \texttt{ShelfStatus} 摊回 Idle；find-then-act 类接口必须原子（如 \texttt{claim\_idle\_by\_role}）。

\begin{lstlisting}[language=Rust, caption=CcAgent::launch\_with\_id 反连握手节选（fuxi-agent-cc/src/agent.rs）, label=lst:cc-launch]
pub async fn launch_with_id(id: AgentId, profile: AgentProfile, cfg: CcLaunchConfig) -> Result<Self> {
    let channel = Arc::new(WsChannel::bind(&Uuid::new_v4().to_string()).await?);
    let mut child = spawn_claude(&cfg)?;          // 注入 --sdk-url + NO_PROXY (避开 Clash TUN)
    tokio::select! {
        Ok(_) = channel.wait_connect(CONNECT_TIMEOUT) => {}
        Ok(_) = child.wait() => return Err(CcError::EarlyExit.into()),
    }
    let pump = tokio::spawn(pump_loop(channel.clone(), /* ... */));
    Ok(CcAgent { card: AgentCard { id, profile, endpoint: channel.url(), .. },
                 inner: Arc::new(Mutex::new(Inner { channel, child, .. })), _pump: pump })
}
\end{lstlisting}

\textbf{三层沙箱 L1/L2/L3。} \texttt{L1ReadOnly} 是项目源码与依赖的只读视图所有门客共享；\texttt{L2Ephemeral} 是任务级临时 worktree，每次 \texttt{dispatch} 创建独立分支、任务结束自动清理；\texttt{L3Persistent} 是角色级持久沙箱，同一角色（如 \texttt{dev-frontend}）跨任务复用同一目录与 build cache。三层中 L2 最关键，落地依赖 git \texttt{worktree}\parencite{git2024worktree}：\texttt{git worktree add -b agent/<id> <path> <base>} 在工作目录之外创建独立工作树与分支，多个 worktree 文件互不影响、共享同一 \texttt{.git} 目录节省磁盘。L3 持久沙箱面向构建缓存复用：编译型语言的 \texttt{target/}、\texttt{node\_modules/} 在每次任务结束时清理代价巨大，L3 让同一角色跨任务复用把 \texttt{cargo build} 从冷启动几分钟降到几秒。L1 只读层面向「门客只读取项目源码」轻量场景，跳过 worktree 创建直接挂载，启动延迟从百毫秒级压到十毫秒级。

\textbf{进程边界 env 清理。} 子进程继承父进程环境变量是 Unix 默认行为，但在 cc/codex 这类「自身就是 cc agent」的 CLI 中会触发嵌套检测——父 cc 起子 cc 时子进程看到 \texttt{CLAUDECODE} 等 env 会静默退出。\texttt{spawn\_claude} 在启动前显式 \texttt{env\_remove} 一组前缀（\texttt{CLAUDECODE}、\texttt{CLAUDE\_CODE\_*}、\texttt{CODEX\_*}），同时注入 \texttt{NO\_PROXY=127.0.0.1,localhost} 规避用户本地 Clash TUN 代理把 WebSocket 反连吞掉的常见环境陷阱。

\subsection{跨节点通讯实现}
\label{subsec:dist}

伏羲 v2 引入跨节点能力：home 机可把任务分发到 worker 节点执行。这套分布式层借鉴 ZooKeeper\parencite{hunt2010zookeeper}「会话 + watcher」心跳模型，但实现做大幅简化——伏羲面向「个人开发者多机部署」不是数据中心规模，没有引入 Raft\parencite{ongaro2014raft} 共识协议，而是接受「home 节点是单点协调者」简化模型。Raft 等协议需至少三节点 quorum 方可保证可用性；个人开发者通常只有一台 mac mini 加一台公司笔记本共两节点，启用 Raft 反因 quorum 不足而不可用。

dist controller 是部署在 home 节点的 axum service，监听同一进程内两套路由：\texttt{/api/*} 走用户 cookie 鉴权（IM 端使用），\texttt{/dist/*} 走 HMAC-SHA256 鉴权（worker 节点使用）。HMAC 签名按全要素 canonical bytes 拼接生成：
\begin{equation}
\sigma = \mathrm{HMAC\text{-}SHA256}\bigl(K,\; \texttt{method} \,\Vert\, \texttt{path} \,\Vert\, \texttt{ts} \,\Vert\, \texttt{nonce} \,\Vert\, \texttt{body}\bigr)
\end{equation}
$K$ 通过 \texttt{FUXI\_DIST\_HMAC\_SECRET} 环境变量注入，五要素覆盖请求方法、路径、毫秒时间戳、随机 nonce 与请求体；验签端按相同公式重算并做常量时间比较（算法~\ref{lst:hmac-verify}）。

\begin{lstlisting}[language=Pseudo, caption=分布式请求 HMAC-SHA256 验签（fuxi-cli/src/dist\_auth.rs）, label=lst:hmac-verify]
1: if |now - ts| > max_skew_ms then return ClockSkew    // 同时防 replay 与 future-replay
2: expected <- HMAC_SHA256(K, method || path || ts || nonce || body)
3: if length(expected) != length(provided_sig) then return BadSignature
4: if not constant_time_eq(expected, provided_sig) then return BadSignature  // 防 timing oracle
5: if not NonceCache.check_and_insert(nonce) then return ReplayedNonce       // 签名验过再查 nonce
\end{lstlisting}

三个安全要点：时间戳偏差检查放最前同时防 replay 与未来时间戳预制；常量时间比较用 \texttt{subtle::ConstantTimeEq::ct\_eq} 防 timing oracle；nonce 检查刻意放在签名验证之后，防止攻击者用大量伪造请求灌满 nonce cache 实现 DoS。worker 节点与 home 之间四条路径：\texttt{/dist/register}（worker 上线声明能力）、\texttt{/dist/heartbeat}（每 2 秒上报心跳与 inflight 计数）、\texttt{/dist/pull}（worker 主动来 home 取活）、\texttt{/dist/event}（worker 把事件回传 home）。设计上保持 stateless：home 不主动推任务给 worker 而是 worker 周期性 pull——这样 worker 在防火墙后也能工作无需暴露公网端口。

worker 节点 inflight 回收靠 \texttt{WorkerHeartbeat} 与 \texttt{WorkerStaleSwept} 两个事件实现。home 端订阅 EventBus 的 watcher 维护 \texttt{Map<NodeId, last\_seen>}，超过 30 秒未收心跳的 worker 被标记 stale，其上 inflight 任务全部 republish 为 \texttt{TaskFailed\{ cause: "worker\_stale" \}}，让玄女决定是否重试。30 秒在 2 秒心跳间隔下相当于允许丢 14 次心跳——过短会让用户在地铁上断网 5 秒就丢任务，过长会让真正死掉的 worker 长时间占着 inflight 槽。worker 端 \texttt{spawn\_worker} 时 \texttt{git fetch} 是 best-effort：失败只 warn 不 fail，避免短暂网络抖动让任务凭空消失。

\subsection{观测、IM 与横向支撑模块}
\label{subsec:firehose-im-support}

\texttt{fuxi-firehose} 内部 \texttt{Hub} 持有 \texttt{EventBus} 引用，对外提供四个端点：\texttt{/ws} 给 TUI 与 CLI 工具消费，\texttt{/sse} 给浏览器与 \texttt{curl} 用户消费，\texttt{/events} 是 REST JSON 历史查询（不做 live tail，回避公理 \#3 反例），还有进程内 broadcast channel 给同进程组件零拷贝消费。同一份事件经四种形态出口、共用底层 \texttt{EventBus::replay} 接口、handler 自身只负责协议翻译——这种「内核统一、出口多样」结构让新增协议（如 gRPC streaming）只需再加一个二三十行 handler。同时支持 WS 与 SSE 两种协议而不强制统一的理由是：TUI 端需要双向（接收事件 + 上报取消信号）WebSocket 自然适配；浏览器与 \texttt{curl} 用户更适合单向 SSE 无需手动管理 ping/pong。

\texttt{fuxi-im} 后端面向 PWA 与桌面 IM 客户端，提供 REST 任务管理（\texttt{GET /api/tasks}、\texttt{POST /api/tasks/<id>/intervene}）、对话历史拉取、双向消息流（\texttt{WS /ws/messages}）与设备推送注册。后端业务调用直接走 \texttt{Arc<Fuxi>} 句柄——\texttt{intervene} 端点直接调 \texttt{Fuxi::intervene}，由编排层判断 idle 退化 dispatch 还是 busy 入 pending 队列，IM 层不重复实现这套语义。\texttt{fuxi-im} 在路由装配上采用「skeleton 稳定、实现可填」分工：\texttt{router.rs} 中所有 handler 签名一次定型，路由装配集中在 \texttt{build()}；具体业务实现按 owner 分布在 \texttt{handlers/*.rs}。

\textbf{横向支撑三件套。} \texttt{fuxi-memory} 基于 SQLite 与 FTS5 全文索引维护三张核心表：\texttt{oracle\_facts}（甲骨：事实三元组）、\texttt{hetu\_patterns}（河图洛书：角色技能样本）、\texttt{user\_profiles}（身份卡）。三张表写入策略统一为「只 ADD 不 UPDATE」——冲突时不覆盖旧行而是把旧行 \texttt{valid\_until} 标 \texttt{now()} 再 INSERT 新行，保证完整审计追溯。FTS5 用 \texttt{unicode61} 分词器，实测 1 万条事实下 p95 查询小于 10 毫秒。\texttt{fuxi-skills} 是角色玉牒加载器，每个角色以 \texttt{roles/<role>/ROLE.md} 存储，借鉴 Voyager\parencite{wang2023voyager} 持久技能库思路并增加 staging + ledger 双轨——新角色先写入 \texttt{roles/<role>.staging/} 经审阅后才 \texttt{mv} 到 \texttt{roles/}。\texttt{fuxi-scheduler} 提供 cron/once/webhook/FileWatch 四类触发器，统一通过发出 \texttt{TriggerFired} 事件接入主链——这条间接路径让 scheduler 自身不需要持有 \texttt{Fuxi} 句柄避免又一处潜在循环依赖。\texttt{fuxi-memory} 的「自动事实抽取」默认关闭，需用户在配置中显式打开——这是因为早期实测发现抽取器在用户与玄女闲聊时容易把无关闲谈写进事实库污染长期记忆；打开后抽取器作为独立子门客被 spawn（进程隔离）从指定 task transcript 中筛候选事实，写库前还会过 \texttt{confidence} 阈值过滤（默认 0.7）。

\subsection{本章小结}

本章把第二章的设计落到代码：模块布局确立 13 crate 边界与「trait 在低层、impl 在高层」反向依赖；事件总线用「\texttt{try\_send} + 后台转交 + 哨兵」三步组合调和实时性与持久化矛盾；A2A 协议从零实现 wire 层、JSON-RPC 信封、SSE 帧解析，并把 \texttt{InputRequired} 中间态在伏羲平台层映射为人工介入信号；编排层用 publish-then-pump 模式把门客私有 receiver 转换为全局事件流，并通过原子 \texttt{claim\_idle\_by\_role} 与 \texttt{shutdown\_agent} 玄女豁免规避并发陷阱；执行层用 WS 反连打破传统 stdio 范式让编排层主动控制门客，三层沙箱让多门客并发文件操作互不干扰；分布式层用 HMAC-SHA256 + nonce cache + 常量时间比较保证跨节点通讯安全；观测与 IM 层把所有事件向多种端形态扇出。下一章在这套实现之上做端到端实测——延迟、吞吐、错误恢复、跨节点协作场景的量化数据。
